%!TEX root = ../main.tex
%!TEX program = xelatex
% Chapter 1: Words 1 - 250
\chapter{The Core Essentials}
\label{chap:chapter1}

{\small\sffamily\color{mutedtext} The most fundamental building blocks of Russian: pronouns, prepositions, and primary everyday verbs.\par}

\vspace{0.4em}
\markboth{Chapter 1: The Core Essentials}{Words 1--250}

\begin{multicols}{2}
\vocentry{1}{И}{i}{And}{Вчера она купила фрукты и овощи.}{She bought some fruit and vegetables yesterday.}
\vocentry{2}{В}{v}{In, on, at}{Маленький котёнок спрятался в подвале.}{A little kitten hid in the basement.}
\vocentry{3}{Не}{nje}{Not, no}{Я не люблю дождливую погоду.}{I do not like rainy weather.}
\vocentry{4}{Он}{on}{He}{Он очень хороший друг.}{He is a very good friend.}
\vocentry{5}{На}{na}{On, at}{Все книги стояли на столе.}{All the books were lying on the table.}
\vocentry{6}{Я}{ja}{I}{Я часто гуляю перед сном.}{I often go for a walk before bedtime.}
\vocentry{7}{Что}{tchto}{That}{Она не знала, что её мать больна.}{She didn't know that her mother was ill.}
\vocentry{8}{Тот}{tot}{That}{Тот парень у окна – мой брат.}{That guy by the window is my brother.}
\vocentry{9}{Быть}{byt’}{To be}{Непросто быть учителем.}{It’s not easy to be a teacher.}
\vocentry{10}{С}{s}{With, from, off, since}{Могу я пойти с тобой?}{May I go with you?}
\vocentry{11}{А}{a}{And, but}{Его младший сын любит спорт, а старший интересуется искусством.}{His younger son likes sports and the elder one is interested in art.}
\vocentry{12}{Весь}{vjes}{All, whole}{Кто съел весь торт?}{Who's eaten all the cake?}
\vocentry{13}{Это}{'ɛtə}{It is, this is}{Это город, в котором выросла моя жена.}{It's the town where my wife grew up.}
\vocentry{14}{Как}{kak}{How, as}{Мы не знаем, как решить эту проблему.}{We don't know how to solve this problem.}
\vocentry{15}{Она}{aˈna}{She}{Я знаю, что она хочет на день рождения.}{I know what she wants for her birthday.}
\vocentry{16}{По}{po}{Along, over, till}{Друзья прогуливались по тихим улицам.}{The friends were walking along the quiet streets.}
\vocentry{17}{Но}{no}{But}{Солнце светит ярко, но на улице холодно.}{The sun is shining brightly but it's cold outside.}
\vocentry{18}{Они}{aˈnji}{They}{Они очень усердно учатся.}{They study very hard.}
\vocentry{19}{К}{k}{To, toward, by}{Не подходи близко к собаке – она может укусить.}{Don’t come close to the dog – it can bite.}
\vocentry{20}{У}{u}{At, by, near}{Не стой у открытого окна, ты можешь простудиться!}{Don’t stand by the open window, you can catch a cold!}
\vocentry{21}{Ты}{ty}{You (2nd person singular)}{Ты слишком поздно ложишься спать.}{You go to bed too late.}
\vocentry{22}{Из}{ɪz}{From, out of}{Я из Лондона.}{I am from London.}
\vocentry{23}{Мы}{my}{We}{Мы так рады видеть вас!}{We are so happy to see you!}
\vocentry{24}{За}{za}{Behind}{Что ты прячешь за спиной?}{What are you hiding behind your back?}
\vocentry{25}{Вы}{vy}{You (2nd person plural or formal singular)}{Дети, куда вы идёте?}{Kids, where are you going?}
\vocentry{26}{Так}{tak}{So, like that}{Я так сильно люблю тебя!}{I love you so much!}
\vocentry{27}{Же}{zhe}{The same (in combination with ‘такой’, ‘тот’, ‘то’)}{Со мной произошла такая же история.}{The same story happened to me.}
\vocentry{28}{От}{ot}{From, since}{Это сообщение от твоей сестры?}{Is this message from your sister?}
\vocentry{29}{Сказать}{skaˈzat’}{To say, to tell}{Я даже не знаю, что сказать.}{I don’t even know what to say.}
\vocentry{30}{Этот}{ˈɛtət}{This, this one}{Тебе нравится этот фильм?}{Do you like this movie?}
\vocentry{31}{Который}{kaˈtoryj}{Which, that, who}{Это тест, который я провалил в прошлом году.}{This is the test, which I failed last year.}
\vocentry{32}{Мочь}{motch}{Сan, be able to (is never used in the initial form)}{Я могу зайти завтра, если хочешь.}{I can drop by tomorrow if you want.}
\vocentry{33}{Человек}{tchela'vek}{A person, a human being, a man}{Её муж очень интересный человек.}{Her husband is a very interesting person.}
\vocentry{34}{О}{o}{About; oh}{Расскажи мне о себе.}{Tell me about yourself.}
\vocentry{35}{Один}{a'din}{One, alone, some}{Мне рассказал об этом один странный человек.}{One strange man told me about it.}
\vocentry{36}{Ещё}{jestch'jo}{Else, more, still}{Что-нибудь ещё?}{Anything else?}
\vocentry{37}{Бы}{by}{A particle used in conditional or subjunctive patterns}{Я бы не доверял им на твоём месте.}{I would not trust them if I were you.}
\vocentry{38}{Такой}{ta'koj}{Such}{Ты такой хороший игрок!}{You're such a good player!}
\vocentry{39}{Только}{ˈtol’kə}{Only, but}{Эти места только для пассажиров с детьми.}{These seats are only for passengers with kids.}
\vocentry{40}{Себя}{seb'ja}{A reflexive pronoun (myself, yourself, herself etc.)}{Она любит себя больше всего на свете.}{She loves herself more than anything in the world.}
\vocentry{41}{Своё}{sva'jo}{A possessive pronoun common for all personal pronouns (my, his, her, its, our, your, their) and refers to thе subject of the sentence. This form is for neuter nouns.}{Он даже не помнит своё имя.}{He doesn’t even remember his name.}
\vocentry{42}{Какой}{ка'koj}{What, which, what kind of}{Какой милый малыш!}{What a cute baby!}
\vocentry{43}{Когда}{kaɡ'da}{When}{Когда ты вернулся?}{When did you come back?}
\vocentry{44}{Уже}{u’zhɛ}{Already}{Мы уже закончили работу.}{We've already finished the work.}
\vocentry{45}{Для}{dlʲa}{For}{Он ничего не сделал для общей победы.}{He didn’t do anything for the common victory.}
\vocentry{46}{Вот}{vot}{Here (is/are)}{Вот ваше место.}{Here is your seat.}
\vocentry{47}{Кто}{kto}{Who}{Люди знают, кто виноват в трагедии.}{People know who is to blame for the tragedy.}
\vocentry{48}{Да}{da}{Yes}{Да, мы уверены, что он прав.}{Yes, we're sure that he is right.}
\vocentry{49}{Говорить}{ɡava'rit’}{To speak, to talk}{Ты мог бы говорить громче, пожалуйста?}{Could you speak up, please?}
\vocentry{50}{Год}{god}{A year}{Это лучший год в моей жизни!}{It's the best year of my life!}
\vocentry{51}{Знать}{znat’}{To know}{Невозможно знать всё.}{It’s impossible to know everything.}
\vocentry{52}{Мой}{moj}{My}{Это мой новый номер телефона, можешь удалить старый.}{This is my new phone number; you can remove the old one.}
\vocentry{53}{До}{do}{Before, till, up to}{До поступления в университет у неё было больше свободного времени.}{Before entering the university, she had more free time.}
\vocentry{54}{Или}{'ili}{Or}{Пара не могла решить, остаться дома или поехать за город.}{The couple couldn’t decide whether to stay at home or to go to the country.}
\vocentry{55}{Если}{'jeslɪ}{If}{Если мы не поторопимся, то опоздаем на поезд.}{If we don't hurry, we'll miss the train.}
\vocentry{56}{Время}{'vremja}{Time}{Время нельзя остановить.}{Time can't be stopped.}
\vocentry{57}{Рука}{ruka}{A hand, an arm}{Мальчик крепко держал маму за руку.}{The boy was firmly holding his mother by the hand.}
\vocentry{58}{Нет}{njet}{No, there is no}{У меня совсем нет времени на спорт.}{I have absolutely no time for sport.}
\vocentry{59}{Самый}{'samyj}{The most}{Это самый захватывающий момент во всей книге.}{It's the most exciting moment in the whole book.}
\vocentry{60}{Ни}{ni}{Neither…nor, not}{Он не умеет играть ни в теннис, ни в гольф.}{He can play neither tennis nor golf.}
\vocentry{61}{Стать}{stat’}{To become, to stand}{Девушка всегда мечтала стать певицей.}{The girl has always dreamt to become a singer.}
\vocentry{62}{Большой}{bal'shoj}{Big, large}{Этот город слишком большой для меня.}{This city is too big for me.}
\vocentry{63}{Даже}{'dazhe}{Even}{Я даже не знаю, как благодарить вас!}{I don’t even know how to thank you!}
\vocentry{64}{Другой}{dru'goj}{Other, another, different}{Вы можете предложить мне какой-нибудь другой цвет?}{Can you offer me any other color?}
\vocentry{65}{Наш}{nash}{Our}{Наш родной язык очень сложный.}{Our mother tongue is very difficult.}
\vocentry{66}{Свой}{svoj}{A possessive pronoun common for all personal pronouns (my, his, her, its, our, your, their) and refers to thе subject of the sentence. This form is for masculine nouns.}{Я не могу вспомнить, где оставил свой зонт.}{I can’t remember where I left my umbrella.}
\vocentry{67}{Ну}{nu}{Well (interjection)}{Ну, как тебе это платье?}{Well, how do you like this dress?}
\vocentry{68}{Под}{pod}{Under, below}{Она нашла потерянную серёжку под кроватью.}{She found the lost earring under the bed.}
\vocentry{69}{Где}{gde}{Where}{Не подскажете, где находится супермаркет?}{Could you tell me where the supermarket is?}
\vocentry{70}{Дело}{d'jelə}{Business, case, matter}{Извини, но это не твоё дело.}{Sorry, but it’s none of your business.}
\vocentry{71}{Есть}{jest’}{To eat, there is/are}{Постарайся не есть после шести.}{Try not to eat after 6 p.m.}
\vocentry{72}{Сам}{sam}{A reflexive pronoun emphasizing the ability to do something without any help, independently. Used for the pronouns ‘He’ and ‘I’ (If ‘I’ is masculine).}{Он может готовить сам, но у него нет на это времени.}{He can cook himself, but he doesn’t have time for it.}
\vocentry{73}{Раз}{raz}{A time}{Не забывай поливать этот цветок хотя бы один раз в неделю.}{Remember to water this flower at least one time a week.}
\vocentry{74}{Чтобы}{ch'toby}{So that, in order to}{Давай составим список покупок, чтобы ничего не забыть.}{Let's make up a shopping list so that we don't forget anything.}
\vocentry{75}{Два}{dva}{Two}{Есть два способа решить эту проблему.}{There are two ways to solve this problem.}
\vocentry{76}{Там}{tam}{There}{Посмотри, что это там такое?}{Look, what is that over there?}
\vocentry{77}{Чем}{tchem}{Than, rather}{Мои новые туфли более удобные, чем старые.}{My old shoes are more comfortable than the old ones.}
\vocentry{78}{Глаз}{glaz}{An eye}{У меня сильно болит правый глаз.}{My right eye hurts badly.}
\vocentry{79}{Жизнь}{zhizn’}{Life}{Она прожила сложную жизнь, но осталась очень доброй.}{She lived a hard life but remained very kind.}
\vocentry{80}{Первый}{'pervyj}{First}{Всегда сложно делать что-то в первый раз.}{It’s always difficult to do something for the first time.}
\vocentry{81}{День}{den’}{Day, daytime}{Врачи рекомендуют начинать день со стакана чистой воды.}{Doctors recommend starting the day with a glass of pure water.}
\vocentry{82}{Тут}{tut}{Here}{Я помню, что оставил ключи тут, но не могу их найти.}{I remember leaving the keys here, but I can’t find them.}
\vocentry{83}{Во}{vo}{In, on, at (a variation of ‘в’ used before consonant clusters)}{Во вторник состоится открытие международной выставки.}{The opening of the international exhibition will take place on Tuesday.}
\vocentry{84}{Ничто}{nich'to}{Nothing}{Ничто не остановит нас на пути к нашей цели.}{Nothing will stop us on the way to our goal.}
\vocentry{85}{Потом}{pa'tom}{Then, afterwards, later}{Идите прямо до перекрестка, а потом поверните налево.}{Go straight ahead up to the crossroads and then turn left.}
\vocentry{86}{Очень}{'otchen’}{Very, very much}{Тренер был очень разочарован результатами команды.}{The coach was very disappointed with the team’s results.}
\vocentry{87}{Со}{so}{With, off, from (a variation of ‘с’ used before consonant clusters)}{Не мог бы ты остаться сегодня со мной?}{Could you stay with me today?}
\vocentry{88}{Хотеть}{ha'tet’}{To want}{Люди могут хотеть разного от жизни.}{People may want different things from life.}
\vocentry{89}{Ли}{li}{Whether, if}{Мне интересно, знают ли они правду.}{I wonder if they know the truth.}
\vocentry{90}{При}{pri}{In the presence of, in the time of}{Не говорите про путешествие при ней – это сюрприз.}{Don’t talk about the trip in her presence, it’s a surprise.}
\vocentry{91}{Голова}{gala'va}{Head}{Моя голова буквально раскалывается от этого шума.}{My head is virtually splitting because of this noise.}
\vocentry{92}{Надо}{'nadə}{Need, must, (it is) necessary}{Мне надо знать, сколько это займёт времени.}{I need to know how much time it will take.}
\vocentry{93}{Без}{bez}{Without}{Я не могу представить себе комнату без окон.}{I can’t imagine a room without windows.}
\vocentry{94}{Видеть}{'videt’}{To see}{Со временем способность видеть ухудшается.}{The ability to see is getting worse with time.}
\vocentry{95}{Идти}{i'ti}{To go, to walk}{Нам лучше идти быстрее, если мы хотим успеть на встречу.}{We’d better go faster if we want to be in time for the meeting.}
\vocentry{96}{Теперь}{te'per’}{Now, in the present}{Теперь, когда ты всё объяснил, ситуация не кажется такой сложной.}{Now that you explained everything the situation doesn’t seem so difficult.}
\vocentry{97}{Тоже}{'tozhe}{Also, too, as well}{Я тоже думаю, что эти перемены к лучшему.}{I also think that these changes are for the better.}
\vocentry{98}{Стоять}{sta'jat’}{To stand}{Нет никакой необходимости стоять в очереди: я уже купила билеты онлайн.}{There is no need to stand in the queue: I have already bought the tickets online.}
\vocentry{99}{Друг}{drug}{A friend}{Настоящий друг всегда рядом.}{A true friend is always there for you.}
\vocentry{100}{Дом}{dom}{A house, home}{Я хочу иметь большой частный дом.}{I want to have a big private house.}
\vocentry{101}{Сейчас}{se'chas}{Now, at the moment}{Сейчас она занята и не может подойти к телефону.}{She is busy now and can’t take the phone.}
\vocentry{102}{Можно}{'mozhnə}{(one) May, (one) can, (it is) possible}{Можно мне воспользоваться твоим ноутбуком?}{May I use your laptop?}
\vocentry{103}{После}{'posle}{After}{Мой дедушка всегда спит после обеда.}{My grandfather always sleeps after dinner.}
\vocentry{104}{Слово}{'slovo}{A word}{Не могли бы вы объяснить мне, что значит это слово?}{Could you explain to me what this word means?}
\vocentry{105}{Здесь}{zdes’}{Here}{Купание здесь запрещено.}{Swimming is forbidden here.}
\vocentry{106}{Думать}{'dumat’}{To think, to suppose}{Я не могу думать, когда вокруг так шумно.}{I can’t think when it’s so noisy around.}
\vocentry{107}{Место}{'mestə}{A place, site}{Это очень мрачное место.}{It’s a very gloomy place.}
\vocentry{108}{Спросить}{spra'sit’}{To ask, to inquire}{Лучше спросить его – он профессионал в этой области.}{It’s better to ask him – he's a professional in the field.}
\vocentry{109}{Через}{'tcherez}{In (time), through, across}{Презентация будет готова через два дня.}{The presentation will be ready in two days.}
\vocentry{110}{Лицо}{li'tso}{A face}{Красивая душа лучше, чем красивое лицо.}{A beautiful soul is better than a beautiful face.}
\vocentry{111}{Что}{tchto}{What}{Я не знаю, что сказать, чтобы приободрить их.}{I don’t know what to say to cheer them up.}
\vocentry{112}{Тогда}{tag'da}{Then, at that time}{Тогда они не знали, что скоро станут богаты.}{They didn’t know then that they would become rich soon.}
\vocentry{113}{Ведь}{ved’}{After all, indeed}{Ведь мы не заставляли его принимать участие!}{After all we didn’t make him participate!}
\vocentry{114}{Хороший}{ha'roshyj}{Good, fine}{Я уверен, это хороший знак.}{I am sure it’s a good sign.}
\vocentry{115}{Каждый}{'kazhdyj}{Every, each}{Он занимается спортом каждый день.}{He practices sport every day.}
\vocentry{116}{Новый}{'novyj}{New}{Власти планируют построить здесь новый район.}{The authorities are planning to build a new district here.}
\vocentry{117}{Жить}{zhyt’}{To live, to reside}{Такие люди не заслуживают того, чтобы жить.}{Such people don’t deserve to live.}
\vocentry{118}{Должный}{'dolzhnyj}{Due, proper}{Очень важно соблюдать должный порядок.}{It’s very important to maintain due order.}
\vocentry{119}{Смотреть}{smat'ret’}{To look, to watch}{Иногда сложно смотреть людям в глаза.}{Sometimes it is difficult to look people in the eye.}
\vocentry{120}{Почему}{patche'mu}{Why}{Я не могу понять, почему ты злишься на меня.}{I can’t understand why you're angry with me.}
\vocentry{121}{Потому}{pata'mu}{Therefore, because of}{Сегодня у меня слишком много работы, потому я не смогу пойти на вечеринку.}{I have too much work today; therefore, I won’t be able to go to the party.}
\vocentry{122}{Сторона}{stara'na}{A side; a party}{Эта сторона улицы более безопасная.}{This side of the street is safer.}
\vocentry{123}{Просто}{p'rostə}{Just; simply}{Я просто хотел убедиться, что всё в порядке.}{I just wanted to make sure that everything is alright.}
\vocentry{124}{Нога}{na'ga}{A leg, a foot}{Он немного хромает, потому что его правая нога короче левой.}{He's a bit lame because his right leg is shorter than the left one.}
\vocentry{125}{Сидеть}{sidet’}{To sit}{Маленькие дети не могут спокойно сидеть.}{Little kids can’t sit still.}
\vocentry{126}{Понять}{pan'jat’}{To understand, to see (perfective)}{Иногда им очень трудно понять друг друга.}{Sometimes it’s very difficult for them to understand each other.}
\vocentry{127}{Иметь}{i'mjet’}{To have}{Чтобы получить эту работу, нужно иметь хорошее образование.}{To get this job one must have a good education.}
\vocentry{128}{Конечный}{ka'nechnyj}{Final}{Руководство интересует конечный результат.}{The management is interested in the final result.}
\vocentry{129}{Делать}{'delat’}{To do; to make}{Я не знаю, что делать в такой ситуации.}{I don’t know what to do in such a situation.}
\vocentry{130}{Вдруг}{vdrug}{Suddenly}{Вдруг тёмная туча закрыла солнце.}{Suddenly a dark cloud covered the sun.}
\vocentry{131}{Над}{nad}{Above, over}{Высоко над крышами быстро летали птицы.}{Birds were quickly flying high above the roofs.}
\vocentry{132}{Взять}{vzjat’}{To take}{Можешь взять все, что тебе нравится.}{You can take anything you like.}
\vocentry{133}{Никто}{nik'to}{Nobody}{Никто не любит плохого отношения.}{Nobody likes bad attitude.}
\vocentry{134}{Сделать}{'zdelat’}{To do, to make (perfective form)}{Что можно сделать, чтобы помочь им?}{What can be done to help them?}
\vocentry{135}{Дверь}{dver’}{A door}{Эта дверь ведёт в большую комнату.}{This door leads to a big room.}
\vocentry{136}{Перед}{'pered}{In front of, before}{Он выступил с речью перед большой аудиторией.}{He made a speech in front of a large audience.}
\vocentry{137}{Нужный}{'nuzhnyj}{Necessary, needed}{Рабочий долго не мог найти нужный инструмент.}{He couldn’t find the necessary tool for a long time.}
\vocentry{138}{Понимать}{pani'mat’}{To understand, to see}{Научись понимать людей вокруг себя.}{Learn to understand people around you.}
\vocentry{139}{Казаться}{ka'zatsa}{To seem}{Молодой человек не хотел казаться слабым.}{The young man didn’t want to seem weak.}
\vocentry{140}{Работа}{ra'bota}{Work; job}{Ведётся работа над восстановлением здания.}{The work on the restoration of the building is being carried out.}
\vocentry{141}{Три}{tri}{Three}{Мне нужно три копии этого документа.}{I need three copies of this document.}
\vocentry{142}{Ваш}{vash}{Your, yours (2nd person plural or 2nd person singular formal)}{Покажите ваш паспорт, пожалуйста.}{Show your passport, please.}
\vocentry{143}{Уж}{uzh}{Emphatic particle}{Я никода не забуду твоих слов, уж это точно!}{I'll never forget your words, that’s for sure!}
\vocentry{144}{Земля}{zem'lja}{Earth; land}{Земля снабжает нас всем необходимым.}{The Earth provides us with everything we need.}
\vocentry{145}{Конец}{ka'nets}{An end; ending}{Конец истории был неожиданным.}{The end of the story was unexpected.}
\vocentry{146}{Несколько}{'neskalkə}{A few, some}{Мне нужно несколько надёжных людей.}{I need a few reliable people.}
\vocentry{147}{Час}{tchas}{An hour}{Остался всего один час до конца представления.}{There's only one hour left before the end of the performance.}
\vocentry{148}{Голос}{'goləs}{A voice}{Мне всегда нравился её мягкий голос.}{I've always liked her soft voice.}
\vocentry{149}{Город}{'gorəd}{A city, a town}{В этот город всегда приятно возвращаться.}{It’s always pleasant to go back to this city.}
\vocentry{150}{Последний}{pas'lednij}{Last, latest}{Последний раз мы виделись три года назад.}{Last time we saw each other three years ago.}
\vocentry{151}{Пока}{pa'ka}{While; for now}{Пока меня не будет, за домом присмотрит моя сестра.}{While I'm away my sister will look after the house.}
\vocentry{152}{Хорошо}{hara'sho}{Well}{Она достаточно хорошо знает русскую литературу.}{She knows Russian literature quite well.}
\vocentry{153}{Давать}{da'vat'}{To give}{Я не могу давать советы в такой ситуации.}{I can’t give advice in such a situation.}
\vocentry{154}{Вода}{va'da}{Water}{Вода – это источник жизни на земле.}{Water is the source of life on earth.}
\vocentry{155}{Более}{'boleje}{More, more than}{Считается, что британцы более консервативны, чем американцы.}{It's considered that the British are more conservative than the Americans.}
\vocentry{156}{Хотя}{ha'tja}{Although, though}{Мы хорошо погуляли, хотя было холодно.}{We had a nice walk, although it was cold.}
\vocentry{157}{Всегда}{vseg'da}{Always}{Я всегда бегаю по утрам.}{I always go jogging in the morning.}
\vocentry{158}{Второй}{vta'roj}{Second}{Второй вариант проекта был более успешным.}{The second variant of the project was more successful.}
\vocentry{159}{Куда}{ku'da}{Where (to)}{Куда мне можно поставить свои вещи?}{Where can I put my things?}
\vocentry{160}{Пойти}{paj'ti}{To go (to)}{Не хочешь ли ты пойти в кино?}{Would you like to go to the cinema?}
\vocentry{161}{Стол}{stol}{A table}{Накрой на стол, пожалуйста.}{Lay the table, please.}
\vocentry{162}{Ребёнок}{reb'jonək}{A child, a kid}{Этот ребёнок очень плохо себя ведёт.}{This child behaves very badly.}
\vocentry{163}{Увидеть}{u'videt’}{To see (perfective)}{Я хочу увидеть Большой Каньон.}{I want to see The Grand Canyon.}
\vocentry{164}{Сила}{'sila}{Power, strength}{Нужна большая сила воли, чтобы бросить курить.}{A strong power of will is needed to give up smoking.}
\vocentry{165}{Отец}{a'tets}{Father}{Отец уделял много времени нашему воспитанию.}{Father devoted a lot of time to our upbringing.}
\vocentry{166}{Женщина}{'zhenschina}{A woman}{Кажется, эта женщина идеальная.}{It seems this woman is perfect.}
\vocentry{167}{Машина}{ma'shina}{A car; machine}{Его новая машина очень дорогая, но надёжная.}{His new car is very expensive but reliable.}
\vocentry{168}{Случай}{'sluchaj}{A case}{Это самый сложный случай в моей карьере.}{It’s the most complicated case in my career.}
\vocentry{169}{Ночь}{notch}{Night}{Ночь была удивительно тёплой и светлой.}{The night was surprisingly warm and light.}
\vocentry{170}{Сразу}{'srazu}{At once}{Я сразу понял, что это ложь.}{I understood at once that it was a lie.}
\vocentry{171}{Мир}{mir}{A world; peace}{Весь мир следит за этими событиями.}{The whole world is following these events.}
\vocentry{172}{Совсем}{sav'sem}{At all}{Мы совсем не возражаем против таких перемен.}{We don’t mind such changes at all.}
\vocentry{173}{Остаться}{as'tatsa}{To stay, to remain}{К сожалению, вам придётся остаться дома.}{Unfortunately, you'll have to stay at home.}
\vocentry{174}{Об}{ob}{About (before vowels)}{Я не могу ничего сказать об остальных актёрах.}{I can’t say anything about the rest of the actors.}
\vocentry{175}{Вид}{vid}{Kind; look; view}{Какой твой любимый вид спорта?}{What is your favorite kind of sport?}
\vocentry{176}{Выйти}{'vyjti}{To go out, to exit}{Наконец, им удалось выйти из лабиринта.}{At last, they managed to go out of the labyrinth.}
\vocentry{177}{Дать}{dat’}{To give (perfective)}{Что они могут дать нам взамен?}{What can they give us in return?}
\vocentry{178}{Работать}{ra'botat’}{To work}{Врачам часто приходится работать сверхурочно.}{Doctors often have to work overtime.}
\vocentry{179}{Любить}{lju'bit’}{To love, to like}{Такие люди неспособны любить.}{Such people aren't able to love.}
\vocentry{180}{Старый}{'staryj}{Old}{Старый рынок популярен среди туристов.}{The old market is popular with tourists.}
\vocentry{181}{Почти}{pach'ti}{Almost}{Они женаты почти два года.}{They've been married for almost two years.}
\vocentry{182}{Ряд}{rjad}{A row; a range}{Билеты в первый ряд уже проданы.}{The tickets in the first row are sold already.}
\vocentry{183}{Оказаться}{aka'zatsa}{To turn out; to find oneself}{Это лекарство может оказаться опасным.}{This medicine can turn out to be dangerous.}
\vocentry{184}{Начало}{na'chalə}{Beginning, start}{Начало пути показалось нам немного трудным.}{The beginning of the way seemed a bit hard to us.}
\vocentry{185}{Твой}{tvoj}{Your (2nd person singular)}{Твой доклад очень подробный.}{Your report is very detailed.}
\vocentry{186}{Вопрос}{vap'ros}{A question; matter}{Можно задать тебе личный вопрос?}{May I ask you a personal question?}
\vocentry{187}{Много}{'mnogə}{Many, much, a lot of}{У меня сегодня слишком много дел.}{I have too many things to do today.}
\vocentry{188}{Война}{vaj'na}{A war}{Любая война – это катастрофа.}{Any war is a disaster.}
\vocentry{189}{Снова}{'snova}{Again}{Когда мы встретимся снова?}{When will we meet again?}
\vocentry{190}{Ответить}{at'vjetit’}{To answer, to reply}{Постарайся ответить на этот вопрос честно.}{Try to answer this question honestly.}
\vocentry{191}{Между}{'mezhdu}{Between}{Какова разница между этими словами?}{What is the difference between these words?}
\vocentry{192}{Подумать}{pa'dumat’}{To think (perfective)}{Мне нужно хорошо подумать, прежде чем я смогу ответить.}{I need to think well before I'll be able to answer.}
\vocentry{193}{Опять}{ap'jat’}{Again}{Он опять сделал ту же самую ошибку.}{He made the same mistake again.}
\vocentry{194}{Белый}{'belyj}{White}{Белый костюм не сочетается с этими зелёными туфлями.}{A white suit doesn't match these green shoes.}
\vocentry{195}{Деньги}{'djengi}{Money}{Нельзя купить счастье за деньги.}{You can’t buy happiness for money.}
\vocentry{196}{Значить}{'znatchit’}{To mean}{Что бы это могло значить?}{What could it mean?}
\vocentry{197}{Про}{pro}{About}{Она и думать не могла про развод.}{She couldn’t even think about a divorce.}
\vocentry{198}{Лишь}{lish’}{Just, only}{Это лишь слова – не принимай близко к сердцу.}{These are just words – don’t take it close to heart.}
\vocentry{199}{Минута}{mi'nuta}{A minute}{Иногда одна минута может длиться очень долго.}{Sometimes one minute can last very long.}
\vocentry{200}{Жена}{'zhena}{A wife}{Его жена прекрасно выглядела на приеме.}{His wife looked wonderful at the reception.}
\vocentry{201}{Посмотреть}{pasmatr'jet’}{To look (perfective); to see}{Мне нужно самой посмотреть на эти файлы.}{I need to look at these files myself.}
\vocentry{202}{Правда}{'pravda}{Truth}{Правда всегда выходит наружу.}{The truth always comes out.}
\vocentry{203}{Главный}{'glavnyj}{Main, chief}{Главный герой пьесы – одинокий инженер.}{The main character of the play is a lonely engineer.}
\vocentry{204}{Страна}{stra'na}{A country}{Вся страна ждёт результатов голосования.}{The whole country is waiting for the results of the vote.}
\vocentry{205}{Свет}{svet}{Light}{Не забудь выключить свет!}{Don’t forget to turn off the light!}
\vocentry{206}{Ждать}{zhdat’}{To wait}{Я просто не могу больше ждать.}{I just can’t wait anymore.}
\vocentry{207}{Мать}{mat’}{A mother}{Мать невесты была строгой женщиной.}{The bride’s mother was a strict woman.}
\vocentry{208}{Будто}{'budtə}{As if; allegedly}{Она так устала, будто работала весь день.}{She was so tired as if she had worked the whole day.}
\vocentry{209}{Никогда}{nikag'da}{Never}{Мы никогда не используем пластиковые пакеты.}{We never use plastic bags.}
\vocentry{210}{Товарищ}{ta'varistch}{Comrade; a friend, a mate}{В СССР люди называли друг друга «товарищ».}{In the USSR people used to call each other ‘comrade’.}
\vocentry{211}{Дорога}{da'roga}{A road, a way}{Эта дорога ведёт к лесу.}{This road leads to the forest.}
\vocentry{212}{Однако}{ad'nakə}{However, yet}{Однако никто не поддержал мою идею.}{However, nobody supported my idea.}
\vocentry{213}{Лежать}{lje'zhat’}{To lie}{Сегодня вечером я хочу лежать на диване и смотреть телевизор.}{Tonight, I want to lie on the sofa and to watch TV.}
\vocentry{214}{Именно}{'imennə}{Exactly, just}{Что именно ты имеешь ввиду?}{What exactly do you mean?}
\vocentry{215}{Окно}{ak'no}{A window}{Закрой окно, пожалуйста, уже холодно.}{Close the window, please, it’s cold already.}
\vocentry{216}{Никакой}{nika'koj}{No, none, not any}{Никакой аргумент не изменит их решение.}{No argument will change their decision.}
\vocentry{217}{Найти}{naj’ti}{To find}{Я не могу найти свои коричневые брюки.}{I can’t find my brown trousers.}
\vocentry{218}{Писать}{pi'sat’}{To write}{Я не умею писать длинные письма.}{I can’t write long letters.}
\vocentry{219}{Комната}{'komnata}{A room}{Детская комната всегда в беспорядке.}{The children’s room is always a mess.}
\vocentry{220}{Москва}{Mosc’va}{Moscow}{Москва – столица Российской Федерации.}{Moscow is the capital of the Russian Federation.}
\vocentry{221}{Часть}{tchast’}{A part}{Это только первая часть концерта.}{It's only the first part of the concert.}
\vocentry{222}{Вообще}{voobs'tche}{In general, generally}{Вообще я не фанат триллеров.}{In general, I'm not a fan of thrillers.}
\vocentry{223}{Книга}{'kniga}{A book}{Книга – лучший подарок.}{A book is the best present.}
\vocentry{224}{Маленький}{'maljenkyj}{Small, little}{Их дом маленький, но очень уютный.}{Their house is small but very cozy.}
\vocentry{225}{Улица}{'ulitsa}{A street}{Эта улица только для пешеходов.}{This street is for pedestrians only.}
\vocentry{226}{Решить}{re'shit'}{To decide}{Компания не может решить, какого кандидата выбрать.}{The company can't decide which candidate to choose.}
\vocentry{227}{Далеко}{dale'ko}{Far away}{Все её родственники живут далеко.}{All her relatives live far away.}
\vocentry{228}{Душа}{du'sha}{A soul}{Душа – самое дорогое, что есть у человека.}{A soul is the most precious thing a person has.}
\vocentry{229}{Чуть}{'tchut'}{A bit, slightly; barely}{Свадебное платье стоило чуть дороже, чем мы ожидали.}{The wedding dress cost a bit more than we expected.}
\vocentry{230}{Вернуться}{ver'nuts'ja}{To return, to come back}{Мы планируем вернуться до полуночи.}{We are planning to return before midnight.}
\vocentry{231}{Утро}{'utrə}{Morning}{То осеннее утро было туманным и холодным.}{That autumn morning was foggy and cold.}
\vocentry{232}{Некоторый}{'nekətəryj}{Some}{В его словах есть некоторый скрытый смысл.}{There's some hidden sense in his words.}
\vocentry{233}{Считать}{schi'tat'}{To consider; to count}{Эту книгу уже можно считать классикой.}{You can already consider this book to be a classic one.}
\vocentry{234}{Сколько}{'skol’kə}{How many, how much}{Сколько гостей вы пригласили на вечеринку?}{How many guests have you invited to the party?}
\vocentry{235}{Помнить}{'pomnit'}{To remember}{Мы всегда будем помнить героизм наших солдат.}{We'll always remember the heroism of our soldiers.}
\vocentry{236}{Вечер}{'vetcher}{Evening}{Мы провели незабываемый романтический вечер.}{We spent an unforgettable romantic evening.}
\vocentry{237}{Пол}{pol}{Floor}{Она уронила кошелёк на пол и не заметила этого.}{She dropped the purse on the floor and didn't notice it.}
\vocentry{238}{Таки}{ta'ki}{Still, in spite of, after all}{Она таки простила своего врага.}{Still, she forgave her enemy.}
\vocentry{239}{Получить}{palu'tchit'}{To get, to receive}{Наш менеджер надеется получить повышение.}{Our manager hopes to get a raise.}
\vocentry{240}{Народ}{na'rod}{People, nation}{Политики часто обманывают народ.}{Politicians often deceive people.}
\vocentry{241}{Плечо}{ple'tcho}{A shoulder}{Он вывихнул плечо во время тренировки.}{He dislocated his shoulder during a workout.}
\vocentry{242}{Хоть}{'hot'}{Though, in spite of}{Хоть она и прожила в Испании десять лет, она не знает языка.}{Though she's been living in Spain for ten years, she doesn't know the language.}
\vocentry{243}{Сегодня}{se'godnja}{Today}{Сегодня мой брат празднует юбилей.}{Today my brother celebrates his anniversary.}
\vocentry{244}{Бог}{bog}{God}{Только Бог может нам помочь!}{Only God can help us!}
\vocentry{245}{Вместе}{v'mestje}{Together}{Это очень дружная семья: они все делают вместе.}{It's a very close-knit family – they do everything together.}
\vocentry{246}{Взгляд}{vzgljad}{View; glance, look}{На мой взгляд, этих денег недостаточно.}{In my view this money is not enough.}
\vocentry{247}{Ходить}{ha'dit'}{To walk, to go}{После операции он снова сможет ходить.}{After the surgery he'll be able to walk again.}
\vocentry{248}{Зачем}{za'tchem}{What for, why}{Зачем тебе нужны эти специи?}{What do you need these spices for?}
\vocentry{249}{Советский}{sa'vjetskij}{Soviet}{Советский Союз был основан в 1922.}{The Soviet Union was founded in 1922.}
\vocentry{250}{Русский}{'ruskij}{Russian}{Русский один из самых сложных языков.}{Russian is one of the most difficult languages.}
\end{multicols}
